% ---
% id: protocol-compiler/rem:exactness
% kind: rem
% title: "Scope of the exact encoding"
% summary: >-
%   The exact choice marginal fixes the sum over outcomes, not the individual
%   probabilities, so no state-independent weight can be factored out of a
%   security argument. Nor does it establish D(iii) of the emitted binary64
%   family, which check V3 decides after the emission decisions.
% prominence: standard
% status: accepted
% uses: [lem:weighted, def:weighted, thm:measurements-and-statistics]
% mentions: [thm:choice-encoding]
% assumes: [A-COIN]
% refs: []
% provenance: newly-derived
% lean:
% ---

\begin{remark}\label{rem:exactness}
Equation \eqref{eq:choicemass} says the choice \emph{marginal} reproduces the
declared weight for every state. What it fixes is that marginal alone: it
does \emph{not} license treating $q_t$ as a fraction that may be factored out
of a security argument, since $q_t$ appears in \eqref{eq:choicemass} only after
summing $\tr[N_{u,t}\omega]$ over $u$, and the terms $\tr[N_{u,t}\omega]$
that carry it depend on $\omega$. That is what separates the emitted encoding
from a fixed fraction applied outside the operators. The exactness statement
matters because it removes the only respectable motive for that unsound
encoding: the weight carried inside the operators already reproduces the
declared choice statistics exactly. Assumption~\ref{ass:coin} is
where the physical content sits --- a choice correlated with the adversary is
not a choice of known bias, and no encoding repairs that.

\emph{Nor does it say that the emitted family passes V3.}
Lemma~\ref{lem:weighted}\,(i) is a statement about the mathematical products
$q_tN_{u|t}$, respectively $P(f\,|\,t(v))F_v$; what is emitted is those
products as represented, as binary64 data, and two things can go wrong between
the two. A product of binary64 numbers need not be binary64, and a family
whose entries are all binary64 may still have a declared-order sum that
rounds. Both are decided before emission, by the emission decisions of
Definition~\ref{def:weighted}, and D(iii) is decided on the emitted blocks
afterwards, by check V3 (Theorem~\ref{thm:measurements-and-statistics}\,(i),
and (i$'$) in the general format), and never inferred from
the lemma, no part of which rests on that theorem. A weighting that fails
either decision is refused, which is the direction that costs a good
declaration rather than admitting a bad one. Both paragraphs are about the products of
Definition~\ref{def:weighted} in either of its forms; what differs between
the forms is only the first: for the setting-dependent form the exactness
statement, Theorem~\ref{thm:choice-encoding}\,(ii), is \emph{not asserted at all}, the
marginal of the choice value being state-dependent through the setting
marginal of the outcome statistics (Lemma~\ref{lem:weighted}\,(iv)), and no
one is tempted to factor it out.
\end{remark}
